\labsection{2}{Files, metadata, and reliable command-line utilities}{sec:lab02}
\begin{labproject}{Inspect a file, copy it with OS interfaces, and build a searchable utility}
\textbf{Suggested time:} 90--110 minutes \hfill
\textbf{Level:} introductory/intermediate\par
\textbf{Starter files:} \texttt{file\_info.c}, \texttt{file\_copy.c}, \texttt{mini\_grep.c}, \texttt{sample.txt}.\par
\textbf{Evidence to keep:} metadata comparison, copy verification, search tests, and exit statuses.

\begin{labmode}
A professional utility must work on normal input and fail clearly on bad input. Every program in this lab is tested both ways.
\end{labmode}

\medskip\noindent\textbf{Learning outcomes.}
\begin{enumerate}
  \item distinguish file content from metadata such as size, permissions, and inode number;
  \item recognise the open/read/write/close resource lifecycle;
  \item use exit status as part of a command-line program's interface.
\end{enumerate}

\medskip\noindent\textbf{Part A --- Inspect metadata.}
\begin{lstlisting}[style=osbash]
cd experiments/Lab02
gcc -Wall -Wextra file_info.c -o file_info
./file_info sample.txt
stat sample.txt
ls -li sample.txt
\end{lstlisting}
Create a copy of the file and compare inode number, size, and permissions. Then rename the original and inspect it again. Explain which properties belong to the file object and which belong to the directory entry/name.

\medskip\noindent\textbf{Part B --- Copy bytes using file descriptors.}
\begin{lstlisting}[style=osbash]
gcc -Wall -Wextra file_copy.c -o file_copy
./file_copy sample.txt copied.txt
cmp sample.txt copied.txt
echo $?
\end{lstlisting}
If \texttt{cmp} prints nothing and returns 0, the files have identical content. Read the source and mark the four stages: acquire descriptors, read, write, release descriptors.

\begin{pitfall}
Do not assume one \texttt{write()} call must always accept every byte requested. The starter code uses a small inner loop so a partial write is handled correctly.
\end{pitfall}

\medskip\noindent\textbf{Part C --- Search and use exit status.}
\begin{lstlisting}[style=osbash]
gcc -Wall -Wextra mini_grep.c -o mini_grep
./mini_grep kernel sample.txt
echo $?
./mini_grep missingword sample.txt
echo $?
./mini_grep kernel does-not-exist.txt
echo $?
\end{lstlisting}
Interpret the three status values. Then use the program in a shell decision:
\begin{lstlisting}[style=osbash]
if ./mini_grep kernel sample.txt >/dev/null; then
    echo "kernel was found"
else
    echo "kernel was not found"
fi
\end{lstlisting}

\medskip\noindent\textbf{Part D --- Make one improvement.}
Choose one of the following: add a byte counter to \texttt{file\_copy}, add a \texttt{--help} message to \texttt{mini\_grep}, or extend \texttt{file\_info} to print whether the path is a directory. Keep the change small and demonstrate it with one test.

\begin{verification}
Submit evidence for one successful and one failing case for each utility. For the copy task, include an independent verification such as \texttt{cmp}; for the search task, include the exit status as well as printed output.
\end{verification}

\begin{tryit}
Optional Linux extension: if \texttt{strace} is installed, run \texttt{strace -e trace=openat,read,write,close ./file\_copy sample.txt copied2.txt}. You do not need to understand every line; identify only the file-related system calls you already know.
\end{tryit}
\end{labproject}
